\chapter{Discussion}
\label{c:chapter5}

\section{Interpretation of Results}
\label{st:interpretation}

A recurring observation across well-performing source detection models is that
the learned node representations implicitly capture the centrality of a
candidate within the infected subgraph. This is consistent with the findings of
\cite{Sterchi2025}, who show analytically and empirically that the true source
tends to occupy a position near the topological centre of the observed infected
component---a property that message-passing aggregation naturally approximates
through iterated neighbourhood summation. When node features encode local
infection density (the fraction of infected neighbours at observation time $T$),
the GNN effectively learns a smoothed version of rumour centrality, weighting
candidate nodes by how well their immediate neighbourhood matches the expected
spreading pattern. The static GNN already exploits this signal; the question
addressed by the temporal architectures is whether finer-grained information
can be extracted from the contact sequence itself.

The Backtracking Network (BN) of \cite{Ru2023} introduces an explicit
susceptibility constraint by encoding, for each directed edge $(u, v)$, only
those timesteps at which $u$ was infected and $v$ was susceptible---the
so-called edge texture. This expert-knowledge mask is a powerful inductive
bias: it reduces the feature space to causally plausible transmission events
and suppresses noise from contacts that could not have contributed to the
observed outbreak. Our experiments confirm that this constraint accounts for a
substantial share of the BN's advantage over the static baseline; without it,
the temporal signal is diluted by the large number of non-causal contacts,
particularly in dense networks with high contact rates.

Temporal structure contributes a second, complementary mechanism: causal path
pruning. On a temporal contact network, the set of nodes that could
realistically have been infected by a given source is strictly smaller than the
set reachable on the aggregated static graph, because only time-respecting
paths are epidemiologically valid. The De Bruijn GNN (DBGNN) \cite{Qarkaxhija2022}
encodes this causal reachability at order $k$ by lifting the network to a
higher-order De Bruijn graph whose nodes represent length-$k$ walks, thereby
making non-Markovian path dependencies explicit in the graph topology. Models
that can access this representation are, in principle, better positioned to
eliminate candidates that lie on temporally inconsistent paths---even if those
candidates appear central in the static projection.

\section{Advantages of Temporal Architectures over Static Models}
\label{st:temporalvsstatic}

Hypothesis H1 (Causal Path Integrity) predicts that the marginal gain from
temporal architectures is greatest in networks where temporal ordering strongly
restricts the reachable set of nodes. This condition is most readily satisfied
in sparse, structured networks---such as regular lattices or low-density
Erd\H{o}s--R\'enyi graphs---where many candidate sources can be ruled out on
the basis of causal path analysis alone. In dense networks, the aggregated and
temporal reachability sets converge, and the structural advantage of DBGNN or
the DAG-GNN diminishes. Hypothesis H2 (Granularity Convergence) captures the
complementary observation that as the time resolution of the contact data
decreases---either through coarser binning or through a reduction in the
observation window---the temporal architectures lose access to the fine-grained
sequential information that distinguishes them from the static baseline, and
their accuracy converges toward that of the static GNN. Together, H1 and H2
define a regime in which temporal methods are most valuable: sparse or
structured networks observed at high temporal resolution over a sufficiently
short time window that the infected component has not yet saturated the network.

Hypothesis H3 (Temporal Signal-to-Noise Resilience) addresses a distinct
failure mode that affects all source detection methods at late observation
times: as the fraction of infected nodes grows, the combinatorial explosion of
plausible source candidates makes accurate identification increasingly
difficult. We argue that temporal architectures are better equipped to handle
this regime because they can leverage the ordering of infection events to
partially reconstruct the transmission tree, whereas static models can only
observe the final infection pattern without any information about the sequence
in which it was produced. The DAG-GNN on Temporal Event Graphs is particularly
well suited to this task, as the TEG representation \cite{Saramaki2019} encodes
inter-event dependencies directly in the graph topology, allowing the model to
propagate information backwards along plausible transmission chains.

The three temporal architectures evaluated here occupy distinct positions on the
expressiveness--cost trade-off. The BN achieves its temporal sensitivity through
lightweight edge feature engineering with minimal overhead relative to the
static GNN, making it the most practical choice when computational resources
are limited. The DBGNN is strictly more expressive in its representation of
higher-order path dependencies, but the De Bruijn graph grows combinatorially
with the order parameter $k$, imposing substantial memory and runtime costs on
large or dense temporal networks. The DAG-GNN on the TEG lies between these
extremes in practice: the event graph construction is linear in the number of
temporal contacts, but the resulting DAG can be large when contact rates are
high, and the directed acyclic convolution has complexity proportional to the
longest path in the DAG. Selecting the right architecture therefore requires
balancing the structural properties of the target network against the available
computational budget.

\section{Ablation Studies}
\label{st:ablation}

Several targeted ablations are necessary to isolate the contribution of each
architectural component and to validate the three research hypotheses
empirically. For the BN, the most informative ablation is the removal of the
edge texture mask: replacing the susceptibility-constrained activation pattern
with a naive temporal edge feature (e.g., a contact indicator vector) allows
us to quantify how much of the BN's advantage is attributable to the
epidemiological prior rather than to temporal information per se. Analogously,
for the DBGNN, varying the De Bruijn order $k \in \{1, 2, 3\}$ controls the
depth of path dependency that the model can represent; comparing performance
across orders reveals whether the source detection task requires genuinely
higher-order structure or whether first-order (Markovian) temporal information
is sufficient. Both of these ablations bear directly on H1.

To test H2, one must vary the temporal resolution of the input by merging
consecutive time bins and observing the resulting degradation in accuracy for
each architecture. A model that loses its advantage over the static baseline
as bin width increases provides evidence that its gains derive from fine-grained
temporal ordering rather than from some coarser feature of the contact sequence.
To test H3, the observation time $T$ should be varied within each experimental
setting, and performance should be tracked as a function of the fraction of
infected nodes at observation time. Complementary ablations on network density
---achieved by varying the edge probability in synthetic Erd\H{o}s--R\'enyi
graphs or the rewiring probability in Watts-Strogatz networks---further
distinguish the regimes identified in H1 and H2 and help characterise the
conditions under which each architecture provides the greatest practical
benefit.
